% !TEX root = ../report.tex

\ssr{ВВЕДЕНИЕ}

Онлайн-маркетплейс объединяет каталог товаров разных продавцов, оформление заказов покупателями, отзывы и административное управление данными. Для такой системы требуется база данных, которая хранит связанные данные о пользователях, товарах, заказах, отзывах и изменении цен, а также поддерживает разграничение доступа между ролями.

Целью курсового проекта является разработка базы данных онлайн-маркетплейса и приложения доступа к ней.

Для достижения поставленной цели необходимо решить следующие задачи:
\begin{itemize}
    \item провести анализ предметной области онлайн-маркетплейса и существующих решений;
    \item сформулировать требования и ограничения к разрабатываемой базе данных и приложению;
    \item выделить роли пользователей и сценарии взаимодействия с системой;
    \item спроектировать сущности базы данных, связи между ними и ограничения целостности;
    \item спроектировать ролевую модель доступа к данным;
    \item спроектировать триггеры журнала изменения цен и пересчёта рейтингов;
    \item спроектировать хранимую функцию расчёта статистики продаж продавца;
    \item выбрать средства реализации базы данных и приложения доступа;
    \item реализовать базу данных, кэширование и интерфейс доступа к данным;
    \item выполнить исследование влияния индексов и Redis-кэширования на измеряемые показатели доступа к данным.
\end{itemize}

\clearpage
